\documentclass[10pt, letterpaper]{article}

% Packages:
\usepackage[
    ignoreheadfoot, % set margins without considering header and footer
    top=2 cm, % seperation between body and page edge from the top
    bottom=2 cm, % seperation between body and page edge from the bottom
    left=2 cm, % seperation between body and page edge from the left
    right=2 cm, % seperation between body and page edge from the right
    footskip=1.0 cm, % seperation between body and footer
    % showframe % for debugging 
]{geometry} % for adjusting page geometry
\usepackage{titlesec} % for customizing section titles
\usepackage{tabularx} % for making tables with fixed width columns
\usepackage{array} % tabularx requires this
\usepackage[dvipsnames]{xcolor} % for coloring text
\definecolor{primaryColor}{RGB}{0, 0, 0} % define primary color
\usepackage{enumitem} % for customizing lists
\usepackage{fontawesome5} % for using icons
\usepackage{amsmath} % for math
\usepackage[
    pdftitle={Magdiel Joshua Angel Galvan's CV},
    pdfauthor={Magdiel Joshua Angel Galvan},
    pdfcreator={LaTeX with RenderCV},
    colorlinks=true,
    urlcolor=primaryColor
]{hyperref} % for links, metadata and bookmarks
\usepackage[pscoord]{eso-pic} % for floating text on the page
\usepackage{calc} % for calculating lengths
\usepackage{bookmark} % for bookmarks
\usepackage{lastpage} % for getting the total number of pages
\usepackage{changepage} % for one column entries (adjustwidth environment)
\usepackage{paracol} % for two and three column entries
\usepackage{ifthen} % for conditional statements
\usepackage{needspace} % for avoiding page brake right after the section title
\usepackage{iftex} % check if engine is pdflatex, xetex or luatex

% Ensure that generate pdf is machine readable/ATS parsable:
\ifPDFTeX
    \input{glyphtounicode}
    \pdfgentounicode=1
    \usepackage[T1]{fontenc}
    \usepackage[utf8]{inputenc}
    \usepackage{lmodern}
\fi

\usepackage{charter}

% Some settings:
\raggedright
\AtBeginEnvironment{adjustwidth}{\partopsep0pt} % remove space before adjustwidth environment
\pagestyle{empty} % no header or footer
\setcounter{secnumdepth}{0} % no section numbering
\setlength{\parindent}{0pt} % no indentation
\setlength{\topskip}{0pt} % no top skip
\setlength{\columnsep}{0.15cm} % set column separation
\pagenumbering{gobble} % no page numbering

\titleformat{\section}{\needspace{4\baselineskip}\bfseries\large}{}{0pt}{}[\vspace{1pt}\titlerule]

\titlespacing{\section}{
    % left space:
    -1pt
}{
    % top space:
    0.3 cm
}{
    % bottom space:
    0.2 cm
} % section title spacing

\renewcommand\labelitemi{$\vcenter{\hbox{\small$\bullet$}}$} % custom bullet points
\newenvironment{highlights}{
    \begin{itemize}[
        topsep=0.10 cm,
        parsep=0.10 cm,
        partopsep=0pt,
        itemsep=0pt,
        leftmargin=0 cm + 10pt
    ]
}{
    \end{itemize}
} % new environment for highlights

\newenvironment{highlightsforbulletentries}{
    \begin{itemize}[
        topsep=0.10 cm,
        parsep=0.10 cm,
        partopsep=0pt,
        itemsep=0pt,
        leftmargin=10pt
    ]
}{
    \end{itemize}
} % new environment for highlights for bullet entries

\newenvironment{onecolentry}{
    \begin{adjustwidth}{
        0 cm + 0.00001 cm
    }{
        0 cm + 0.00001 cm
    }
}{
    \end{adjustwidth}
} % new environment for one column entries

\newenvironment{twocolentry}[2][]{
    \onecolentry
    \def\secondColumn{#2}
    \setcolumnwidth{\fill, 4.5 cm}
    \begin{paracol}{2}
}{
    \switchcolumn \raggedleft \secondColumn
    \end{paracol}
    \endonecolentry
} % new environment for two column entries

\newenvironment{threecolentry}[3][]{
    \onecolentry
    \def\thirdColumn{#3}
    \setcolumnwidth{, \fill, 4.5 cm}
    \begin{paracol}{3}
    {\raggedright #2} \switchcolumn
}{
    \switchcolumn \raggedleft \thirdColumn
    \end{paracol}
    \endonecolentry
} % new environment for three column entries

\newenvironment{header}{
    \setlength{\topsep}{0pt}\par\kern\topsep\centering\linespread{1.5}
}{
    \par\kern\topsep
} % new environment for the header

\newcommand{\placelastupdatedtext}{% \placetextbox{<horizontal pos>}{<vertical pos>}{<stuff>}
  \AddToShipoutPictureFG*{% Add <stuff> to current page foreground
    \put(
        \LenToUnit{\paperwidth-2 cm-0 cm+0.05cm},
        \LenToUnit{\paperheight-1.0 cm}
    ){\vtop{{\null}\makebox[0pt][c]{
        \small\color{gray}\textit{Last updated in September 2024}\hspace{\widthof{Last updated in September 2024}}
    }}}%
  }%
}%

% save the original href command in a new command:
\let\hrefWithoutArrow\href

% new command for external links:

\begin{document}
    \newcommand{\AND}{\unskip
        \cleaders\copy\ANDbox\hskip\wd\ANDbox
        \ignorespaces
    }
    \newsavebox\ANDbox
    \sbox\ANDbox{$|$}

    \begin{header}
        \fontsize{25 pt}{25 pt}\selectfont Magdiel Joshua Angel Galvan

        \vspace{5 pt}

        \normalsize
        \mbox{Cuautitlan Izcalli, Mex}%
        \kern 5.0 pt%
        \AND%
        \kern 5.0 pt%
        \mbox{\hrefWithoutArrow{mailto:magdiel50.mag@gmail.com}{magdiel50.mag@gmail.com}}%
        \kern 5.0 pt%
        \AND%
        \kern 5.0 pt%
        \mbox{\hrefWithoutArrow{tel:+525611130971}{+52 5611130971}}%
        \kern 5.0 pt%
        \AND%
        \kern 5.0 pt%
        \mbox{\hrefWithoutArrow{https://www.linkedin.com/in/magdiel-joshua-angel-galvan-991664345}{LinkedIn}}%
        \kern 5.0 pt%
        \AND%
        \kern 5.0 pt%
        \mbox{\hrefWithoutArrow{https://github.com/magdielzz}{GitHub}}%
    \end{header}

    \vspace{5 pt - 0.3 cm}

    \section{About me}
        
        \begin{onecolentry}
     Fisico de la UNAM especializado en ciencia de datos y modelado computacional. Actualmente me desempeño como Business Analyst en Orion Innovation, colaborando con El Palacio de Hierro en proyectos de análisis de datos y Business Intelligence: análisis exploratorio de datos (EDA), modelos analíticos de clustering, forecasting y análisis de curvas, tableros interactivos en Tableau, y automatización de reportes mediante repositorios en Git desplegados a Looker, apoyándome en Python, SQL y AWS. Cuento con experiencia técnica en investigación de Machine Learning (PINNs y CNNs) y análisis estadístico, enfocándome en resolver problemas complejos mediante el desarrollo de software, la computación paralela y la visualización de datos.
   
        \end{onecolentry}

    \section{Education}
        
        \begin{twocolentry}{
            2022 -- 2026
        }
            \textbf{Facultad de Ciencias, Universidad Nacional Autónoma de México (UNAM), CDMX, México}, Licenciatura en Física (Titulado)\end{twocolentry}

        \vspace{0.10 cm}
        \begin{onecolentry}
            \begin{highlights}
                \item Enfasis en pensamiento analítico y matemático aplicado a problemas científicos y técnicos.
                \item \textbf{Cursos relevantes:} Física Computacional, Matemáticas Avanzadas, Ciencia de Datos.
                \item \textbf{Tesis:} \textit{Estudio de redes neuronales informadas por la física para la dinámica multiperiódica de largo tiempo en ecuaciones diferenciales no lineales.}
            \end{highlights}
        \end{onecolentry}
        
        \begin{twocolentry}{
            2019 -- 2022
        }
            \textbf{Colegio de Ciencias y Humanidades Plantel Azcapotzalco}, Bachillerato y Técnico en Desarrollo de Software y Sistemas Computacionales\end{twocolentry}

        \vspace{0.10 cm}
        \begin{onecolentry}
            \begin{highlights}
                \item Formación en programación y desarrollo de software con enfoque práctico.
                \item Proyectos prácticos en desarrollo de aplicaciones y sistemas.
            \end{highlights}
        \end{onecolentry}

    \section{Experience}
    \begin{twocolentry}{
    2026 -- Actualidad
}

\textbf{Business Specialist Analyst Data}, Orion Innovation (cuenta El Palacio de Hierro)
\end{twocolentry}
\vspace{0.10 cm}
\begin{onecolentry}
    \begin{highlights}
        \item Colaboración con El Palacio de Hierro en proyectos de análisis de datos y Business Intelligence, realizando análisis exploratorio de datos (EDA) para identificar patrones de consumo y oportunidades de negocio.

        \item Desarrollo y optimización de reportes mediante la creación de repositorios en Git, desde donde se despliegan automáticamente a Looker para la generación de reportes recurrentes sin intervención manual.

        \item Construcción de modelos analíticos aplicando técnicas de clustering, forecasting y análisis de curvas para apoyar la toma de decisiones estratégicas del negocio.

        \item Diseño y mantenimiento de tableros interactivos en Tableau para el seguimiento de indicadores clave y la visualización de resultados ante distintas áreas.

        \item Automatización de procesos de análisis y visualización de datos utilizando Python, Git y servicios de AWS, transformando datos crudos en información accionable para la toma de decisiones.

        \item Desarrollo de automatizaciones mediante Docker y Google Cloud Run, mejorando la ejecución y escalabilidad de procesos operativos y analíticos.

        \item Participación en la elaboración de un agente de accionables para el área de Mercadeo, orientado a convertir información analítica en recomendaciones para la toma de decisiones.

        \item Coordinación de actividades y planeación de procesos de trabajo en proyectos de las áreas de Contenido y Comercial, liderando tanto la definición intelectual de los procesos como su ejecución con equipos multidisciplinarios.

        \item Desarrollo de nuevos tableros para dar visibilidad al desempeño de productos en distintos puntos de venta, incluyendo tiendas físicas y e-commerce.

        \item Perfeccionamiento de un modelo de asignación de productos a tiendas que considera múltiples factores para maximizar ventas, reducir inventario estancado y mitigar faltantes de SKU en puntos de venta con mayor demanda.
    \end{highlights}
\end{onecolentry}
        \begin{twocolentry}{
            2025 -- 2026
        }           
    \textbf{Data Analyst}, Publicis Groupe
\end{twocolentry}
\vspace{0.10 cm}
\begin{onecolentry}
    \begin{highlights}
        \item Elaboración y análisis de reportes estratégicos para Ferrero, evaluando el desempeño de marcas de chocolates, desayunos y productos infantiles mediante métricas de redes sociales, tendencias de consumo y comparativos históricos para identificar oportunidades de optimización en campañas de marketing.

        \item Desarrollo de soluciones de automatización en Python (Pandas) y Excel para la actualización de reportes de Banco Azteca, integrando nuevas bases de datos provenientes de equipos externos y optimizando los procesos de limpieza, transformación y actualización de información.

        \item Análisis de inversión publicitaria, impresiones, creativos, medios, cobertura geográfica y estrategias de competidores para Banco Azteca, generando recomendaciones basadas en datos para optimizar la distribución del presupuesto publicitario y maximizar el alcance de campañas.

        \item Construcción de dashboards y reportes ejecutivos en Excel y Power BI para FarmaMédica, integrando información proveniente de HR Media y otras fuentes para el seguimiento de indicadores clave de desempeño (KPIs).

        \item Desarrollo de estudios de audiencias mediante GWI, TGI, Epsilon ID Modeler y técnicas de clustering con Scikit-learn para la segmentación de consumidores y la generación de perfiles estratégicos utilizados en campañas de marketing.

        \item Elaboración de investigaciones de mercado utilizando plataformas como Pulsar, Euromonitor, Admetricks y Auditsa para analizar tendencias de consumo, Social Listening, publicidad exterior (OOH), comportamiento de categorías y estrategias de la competencia.

        \item Automatización de procesos ETL, extracción, limpieza y transformación de datos mediante Python, Excel y consultas básicas en SQL, estandarizando flujos de trabajo y mejorando la eficiencia del equipo de Business Intelligence.

        \item Presentación y comunicación de resultados analíticos a equipos internos y clientes, exponiendo insights estratégicos y recomendaciones basadas en datos para la toma de decisiones en campañas de medios y marketing digital.
    \end{highlights}
\end{onecolentry}
        \begin{twocolentry}{
            2025
        }
            \textbf{Colaborador en Proyectos Emprendedores}, Seminario de emprendimiento, Facultad de ciencias CA\end{twocolentry}

        \vspace{0.10 cm}
        \begin{onecolentry}
            \begin{highlights}
                \item Participé en el desarrollo de un proyecto tecnológico en etapas iniciales, colaborando en un equipo multidisciplinario.
                \item Implementé soluciones avanzadas de ciencia de datos utilizando Python para automatización, análisis predictivo y modelado estadístico.
                \item Desarrollé e implementé modelos de machine learning incluyendo Árboles de Decisión, Bosques Aleatorios y XGBoost para problemas de clasificación y regresión.
                \item Construí Redes Neuronales con TensorFlow/Keras, incluyendo un modelo para predicción de tumores malignos a partir de imágenes médicas.
                \item Apliqué algoritmos no supervisados como Agrupación K-Means para segmentación de datos y K-Vecinos Más Cercano para problemas de clasificación.
                \item Desarrollé habilidades en metodologías ágiles y resolución de problemas en entornos dinámicos.
            \end{highlights}
        \end{onecolentry}

        \vspace{0.2 cm}

        \begin{twocolentry}{
            2023--2025
        }
            \textbf{Aplicación de Física Computacional}, Facultad de ciencias UNAM\end{twocolentry}

        \vspace{0.10 cm}
        \begin{onecolentry}
            \begin{highlights}
                \item Desarrollé un modelo computacional en Python para simular fenómenos físicos.
                \item Desarrollé un solucionador numérico de las ecuaciones de Navier-Stokes en formulación vorticidad-escalar, implementando métodos pseudo-espectral en 2D y 3D.
                \item Implementé el algoritmo en Python utilizando NumPy, SciPy y TensorFlow.
                \item Optimicé el cómputo paralelizando las transformadas de Fourier.
            \end{highlights}
        \end{onecolentry}
        \begin{twocolentry}{
            2025 -- Actualidad
        }
            \textbf{Investigación y Desarrollo de PINNs}, Facultad de Ciencias UNAM\end{twocolentry}

        \vspace{0.10 cm}
        \begin{onecolentry}
            \begin{highlights}
                \item Implementación de Redes Neuronales Informadas por la Física (PINNs) para la resolución de ecuaciones diferenciales parciales complejas.
                \item Ejecución de simulaciones y entrenamiento de modelos en el clúster del Instituto de Ciencias Nucleares, utilizando Linux y herramientas de cómputo de alto rendimiento.
            \end{highlights}
        \end{onecolentry}

        \vspace{0.2 cm}

        \begin{twocolentry}{
            2025
        }
            \textbf{Clasificación de Galaxias mediante Deep Learning}, Facultad de Ciencias UNAM\end{twocolentry}

        \vspace{0.10 cm}
        \begin{onecolentry}
            \begin{highlights}
                \item Desarrollo de modelos de clasificación para galaxias con y sin estructuras de marea utilizando visión computacional.
            \end{highlights}
        \end{onecolentry}

        \vspace{0.2 cm}

        
\begin{twocolentry}{Marzo 2026}
    \textbf{Presentador de Póster}, 2° Congreso de Ciencia, Tecnología y Vinculación: Innovación para la Salud 2026
\end{twocolentry}
\vspace{0.10 cm}
\begin{onecolentry}
    \begin{highlights}
        \item Presentación de póster de investigación.
    \end{highlights}
\end{onecolentry}
\begin{twocolentry}{Mayo 2026}
    \textbf{Presentador Galardonado de Póster}, XXXIV Reunión Anual de la DGFM - Sociedad Mexicana de Física
\end{twocolentry}
\vspace{0.10 cm}
\begin{onecolentry}
    \begin{highlights}
        \item Presentación de póster sobre la dinámica multi-periódica en ecuaciones diferenciales no lineales mediante redes neuronales informadas por la física (PINNs).
        \item Distinción en el evento tras calificar dentro de la selección de mejores pósters de investigación.
    \end{highlights}
\end{onecolentry}
\begin{twocolentry}{2025 -- 2026}
    \textbf{Ayudante de Profesor de Ciencia de Datos}, UNAM
\end{twocolentry}
\vspace{0.10 cm}
\begin{onecolentry}
    \begin{highlights}
        \item Ayudante durante dos semestres consecutivos (un año) en cursos de Ciencia de Datos e Inteligencia Artificial para grupos presenciales de aproximadamente 40 estudiantes.
        \item Asistencia técnica en la programación y depuración de código en Python utilizando conjuntos de datos reales y sintéticos.
        \item Apoyo en la enseñanza teórica y práctica de algoritmos esenciales de Machine Learning como clustering, regresión lineal/logística y modelos de clasificación.
    \end{highlights}
\end{onecolentry}
    \section{Languages}
        
        \begin{samepage}
            \begin{twocolentry}{
                
            }
                \textbf{Español (nativo)}
            \end{twocolentry}

            \vspace{0.10 cm}
            \begin{onecolentry}
                \textbf{Inglés (Intermedio)} 
            \end{onecolentry}
        \end{samepage}

    \section{Projects}

        \begin{twocolentry}{
            Tesis de licenciatura
        }
            \textbf{Estudio de redes neuronales informadas por la física para la dinámica multiperiódica de largo tiempo en ecuaciones diferenciales no lineales}\end{twocolentry}

        \vspace{0.10 cm}
        \begin{onecolentry}
            \begin{highlights}
                \item Investigación y desarrollo de Redes Neuronales Informadas por la Física (PINNs) para estudiar la dinámica de largo tiempo en ecuaciones diferenciales no lineales.
                \item Uso del clúster del Instituto de Ciencias Nucleares en Linux para ejecutar experimentos computacionales y entrenamiento de modelos en un entorno de cómputo de alto rendimiento.
            \end{highlights}
        \end{onecolentry}

        \vspace{0.2 cm}
        
        \begin{twocolentry}{
            \href{https://github.com/magdielzz/Navier-Stokes-Parallel-Solver.git}{github.com/magdielzz/Navier-Stokes-Parallel-Solver}
        }
            \textbf{Navier Stokes Parallel Solver}\end{twocolentry}

        \vspace{0.10 cm}
        \begin{onecolentry}
            \begin{highlights}
                \item Este proyecto implementa un solucionador para las ecuaciones de Navier-Stokes en 2D en forma de vorticidad usando un método pseudo-espectral utilizando un enfoque de computación paralela.
                \item Tools Used: Python
            \end{highlights}
        \end{onecolentry}

        \vspace{0.2 cm}

        \begin{twocolentry}{
            \href{https://github.com/magdielzz/Proyecto-ciencia-de-Datos-Biomedicina.git}{github.com/magdielzz/Proyecto-ciencia-de-Datos-Biomedicina}
        }
            \textbf{Análisis de Estancias Hospitalarias Prolongadas para México con Ciencia de Datos}\end{twocolentry}

        \vspace{0.10 cm}
        \begin{onecolentry}
            \begin{highlights}
                \item El proyecto analiza la duración de la estancia hospitalaria (LOS) como indicador de retrasos en el diagnóstico y tratamiento de enfermedades en México.
                \item Tools Used: Python
            \end{highlights}
        \end{onecolentry}

        \vspace{0.2 cm}
        \begin{twocolentry}{
            \href{https://github.com/magdielzz/An-lisis-trastornos-sue-o.git}{github.com/magdielzz/An-lisis-trastornos-sue-o.git}
        }
            \textbf{Análisis de Trastornos del Sueño}\end{twocolentry}

        \vspace{0.10 cm}
        \begin{onecolentry}
            \begin{highlights}
                \item Este repositorio contiene un análisis integral de trastornos del sueño utilizando técnicas de Machine Learning.
                \item Tools Used: Python
            \end{highlights}
        \end{onecolentry}
        \begin{twocolentry}{
            \href{https://github.com/magdielzz/PINNs}{github.com/magdielzz/PINNs}
        }
            \textbf{Physics-Informed Neural Networks (PINNs)}\end{twocolentry}

        \vspace{0.10 cm}
        \begin{onecolentry}
            \begin{highlights}
                \item Implementación de redes neuronales que incorporan leyes físicas para modelado avanzado.
            \end{highlights}
        \end{onecolentry}

        \vspace{0.2 cm}

        \begin{twocolentry}{
            \href{https://github.com/magdielzz/Clasificaci-n-de-galaxias-con-estructura-de-marea-y-sin-estructura-de-marea}{github.com/magdielzz/Galaxy-Classification}
        }
            \textbf{Clasificación de Galaxias}\end{twocolentry}

        \vspace{0.10 cm}
        \begin{onecolentry}
            \begin{highlights}
                \item Uso de Redes Neuronales Convolucionales para la identificación de estructuras astronómicas.
            \end{highlights}
        \end{onecolentry}

    \section{Skills}
        
        \begin{onecolentry}
            \textbf{Languages and Tools:} C, Python (Pandas, NumPy, SciPy, TensorFlow/Keras, Matplotlib), Java, R, SQL, Docker, Google Cloud Run, GitHub, Tableau, Power BI, Excel Avanzado, GWI, TGI, Epsilon ID Modeler, Pulsar, HR Media, Euromonitor, Admetricks, Auditsa, Git, AWS, Looker, Linux y herramientas de cómputo de alto rendimiento (clúster).
        \end{onecolentry}

        \vspace{0.2 cm}

        \begin{onecolentry}
            \textbf{Competencias Analíticas:} Resolución de problemas, pensamiento matemático, modelado computacional, análisis exploratorio de datos, ciencia de datos, conceptos de inteligencia artificial generativa, aplicaciones de SQL en entornos reales.
            
        \end{onecolentry}

        \vspace{0.2 cm}

        \begin{onecolentry}
            \textbf{Habilidades Interpersonales:} Trabajo en equipo multidisciplinario, adaptabilidad, comunicación efectiva, 
        \end{onecolentry}

\end{document}
